\section{Discussion}\label{sec:discussion}
RQ1 is answered positively within the observed boundary: the public MV name moved from baseline to candidate logic without a query-loop error. The runbook-recorded duration of \texttt{0.145226 s} is an operational indicator for a small SQL-logic change on this laboratory cluster; since file provenance does not enable independent latency reconstruction, it is neither an SLA nor proof for all query patterns or failure modes.

RQ2 is supported by Parquet and metadata objects in MinIO after replay. This demonstrates an operational sink in the exercised pipeline; because RisingWave documents the swap as an MV name operation and treats dependent sinks separately, the study does not claim that the sink moved to green logic after the swap.

RQ3 is supported by cluster health and resource dashboards. The low rate limit and the deliberate shutdown of Vector protect the resource-constrained data-plane node; the result establishes functional feasibility, not maximum capacity.

Relative to Ursa, Continux does not propose a streaming engine. It contributes an operational perspective that combines lakehouse-streaming motivation with GitOps, observability, and public query-logic replacement on a modest deployed environment.

\subsection{Limitations}\label{sec:limitations}
\begin{table*}[t!]
\centering
\caption{Evidence limitations and consequences}
\begin{tabularx}{\textwidth}{@{}p{0.49\textwidth}Y@{}}
\toprule
\textbf{Limitation} & \textbf{Consequence}\\
\midrule
No offset, stable-event-key, duplicate, or loss reconciliation & End-to-end exactly-once is not verified\\
Replay stopped at \texttt{986} trips & No full-dataset or peak-throughput claim\\
No measured runs for three load levels & No load/latency benchmark is reported\\
Sink dependency predates MV swap & No Iceberg green-cutover claim\\
Incomplete Kafka catalog/Iceberg freshness metrics & Findings rely on SQL counts and MinIO listings\\
Duration/swap files updated after the transition log & Duration is reported as a runbook measure, not independently reconstructed latency\\
Topic replication factor is \texttt{1} & Broker node-failure tolerance is not evaluated\\
Raw dataset can contain \texttt{pickup\_time} outside the configured month and extreme \texttt{trip\_distance} outliers & Current green logic filters only negative values; time-range and distance-range filters are not applied\\
\bottomrule
\end{tabularx}
\end{table*}

\subsection{Future work}
Future research can pursue four complementary directions. First, attach a stable \emph{event identity} -- an identifier preserved from the source JSONL through to the Iceberg sink -- so that duplicates and losses can be reconciled across the entire path, allowing end-to-end exactly-once to be verified by measurement rather than inferred from connector configuration. Second, run measured multi-load benchmarks on a larger data plane, using the load profiles defined in the repository (\texttt{low}, \texttt{medium}, \texttt{high}) to produce latency/throughput curves. Third, design a cutover that also covers the downstream sink -- for instance creating a new sink bound to the green MV before cutover and renaming the sink, or compacting and rewriting Iceberg snapshots under the new logic -- so that conclusions can be extended to the full lakehouse path. Fourth, extend the source from a recorded dataset to live mobility events (for example traffic sensors or bus GPS telemetry). These are directions for subsequent releases, not part of the current results.
